% ColorSAR - IEEE conference paper, DRAFT 1 (18 Nov 2026).
% Build on Overleaf: upload colorsar.tex + refs.bib, compiler pdfLaTeX, run BibTeX.
%
% Every number in this draft comes from a run log on the 7 held-out test scenes (n = 1660):
%   all_results/picma/results/{histmatch,welsh,unet_r50}.log   (Kaggle colorsar-bench v2)
%   all_results/picma_t4/results/picma.log                     (Kaggle colorsar-bench v3, 1 h budget)
%   Priyanshaur/ColorSAR @ diffusion-experiment b5cc366, diffusion_small_run/metrics_test.json
% Red \pend{...} text = waiting for a run or a section owner. Do not invent numbers for it.
% Section owners follow docs/WORK_SPLIT.md.

\documentclass[conference]{IEEEtran}
\IEEEoverridecommandlockouts
\usepackage{cite}
\usepackage{amsmath,amssymb,amsfonts}
\usepackage{graphicx}
\usepackage{booktabs}
\usepackage{multirow}
\usepackage{xcolor}
\usepackage{url}
\newcommand{\pend}[1]{\textcolor{red}{[#1]}}

\begin{document}

\title{ColorSAR: Physics-Informed Cross-Modal Attention Latent Diffusion\\
for SAR-to-Optical Image Translation}

\author{
\IEEEauthorblockN{\pend{Author 1}, \pend{Author 2}, \pend{Author 3}, \pend{Author 4}, \pend{Author 5}}
\IEEEauthorblockA{\pend{Department, Institute, City, Country}\\
\pend{emails}}
}

\maketitle

% Owner: Person 4 - rewrite last, once all rows are in.
\begin{abstract}
Synthetic aperture radar (SAR) images the Earth through clouds and at night, but its grey,
speckled images are hard for people to read. We study the translation of a Sentinel-1 SAR
tile into the matching Sentinel-2 optical tile. We propose PICMA, a physics-informed
cross-modal attention latent diffusion model. PICMA feeds the generator a three-channel
physics input built with the Lee speckle filter and adds a speckle-aware physics loss. It
aligns the SAR encoder with an optical encoder by contrastive pretraining, and it conditions
a latent denoiser on 320 SAR tokens through cross-attention. We compare PICMA with
\pend{eight} methods, from classical colour transfer to GANs, Transformers and pixel-space
diffusion. Every method uses the same scene-level split, so the 1{,}660 test tiles come from
seven Sentinel scenes that no model saw in training. All methods are scored by the same code
in CIE Lab and with perceptual metrics. On this strict split, no method beats a constant
mean colour on a/b colour error. Classical statistics transfer is best on PSNR and LPIPS. PICMA, trained
for only one hour on two T4 GPUs, has the lowest FID (244.9), but by a margin within the
noise of the metric, and it is the worst on the pixel metrics.
\pend{update with the remaining baselines and the ablation}.
\end{abstract}

\begin{IEEEkeywords}
SAR-to-optical translation, Sentinel-1, Sentinel-2, speckle, Lee filter, latent diffusion,
cross-attention, contrastive learning, image colorization
\end{IEEEkeywords}

% =====================================================================================
% Owner: Person 4
\section{Introduction}

Optical satellites such as Sentinel-2 record colour images of the Earth's surface, but they
cannot see through clouds and they need daylight. In monsoon regions, an area can stay
hidden from optical sensors for weeks. SAR satellites such as Sentinel-1 send microwave pulses
and record the echo. They work through clouds, day and night. A SAR image, however, records
surface roughness and geometry, not colour. It is grey and covered in multiplicative speckle
noise, so it is hard for a non-expert to read.

SAR-to-optical translation predicts the optical image that a cloud-free optical sensor would
have recorded. The task is ill-posed. Two fields with the same radar texture can be green or
brown, so one SAR tile matches many plausible colourings. A network trained with a pixel loss
therefore learns to predict the average colour. Its output is dull and greyish.
Generative models such as GANs and diffusion models can sample one sharp, plausible
colouring instead. They then trade pixel accuracy for realism, which is the
perception-distortion trade-off~\cite{blau2018perception}.

A second problem is evaluation. Tiles cut from one satellite scene overlap and look alike.
When train and test tiles are drawn at random from the same scenes, the test score rewards
memorisation. In our own earlier experiments, a per-pair random split inflated
the test PSNR of the SAR colourizer to 16.6~dB. We therefore hold out whole scenes.

\textbf{Contributions.}
\begin{itemize}
\item PICMA, a latent diffusion model for SAR-to-optical translation. It combines (i) a
  physics input and a physics loss derived from the multiplicative speckle model and the
  Lee filter~\cite{lee1980speckle}, (ii) contrastive SAR--optical pretraining of the
  conditioning encoder, (iii) cross-attention from the denoiser to SAR tokens, and (iv) a
  focal frequency loss~\cite{jiang2021ffl} on the decoded image.
\item A benchmark of nine methods on a scene-held-out split of Sentinel-1/2 pairs. All methods
  share the data, the 1{,}660 test tiles and the scoring code, and the training budget is
  reported for each.
\item An honest analysis: on this split, the learned models do not beat a constant mean
  colour on a/b error, and we discuss why.
\end{itemize}

% =====================================================================================
% Owner: Person 3 - extend to 15-20 papers, one paragraph per baseline family.
\section{Related Work}

\textbf{Classical colourization.} Welsh et al.~\cite{welsh2002transferring} transfer colour
from a reference image to a grey image by matching local luminance statistics. It needs no
training. We use a SAR-adapted version as a baseline.

\textbf{CNN colourization.} Zhang et al.~\cite{zhang2016colorful} cast colourization as
classification over quantised a/b bins with class rebalancing, and decode with an annealed
mean. Our U-Net baseline uses this head on a ResNet-50~\cite{he2016resnet}
U-Net~\cite{ronneberger2015unet}.

\textbf{GANs.} Pix2Pix~\cite{isola2017pix2pix} trains a U-Net generator against a PatchGAN
discriminator on paired data. CycleGAN~\cite{zhu2017cyclegan} drops the pairing with a
cycle-consistency loss. Fuentes Reyes et al.~\cite{fuentesreyes2019sar2opt} studied
conditional GANs for SAR-to-optical translation on SEN1-2~\cite{schmitt2018sen12} and
discussed the limits of the task.

\textbf{Transformers and frequency methods.} TransUNet~\cite{chen2021transunet} places a
Transformer~\cite{vaswani2017attention} in the U-Net bottleneck for global context.
Multi-level wavelet CNNs~\cite{liu2018mwcnn} replace pooling with the discrete wavelet
transform. The focal frequency loss~\cite{jiang2021ffl} weights frequency errors by their size
and so reduces blur.

\textbf{Diffusion models.} DDPM~\cite{ho2020ddpm} learns to reverse a gradual noising process.
DDIM~\cite{song2021ddim} samples deterministically in far fewer steps.
Palette~\cite{saharia2022palette} applies conditional diffusion to image-to-image tasks,
colourization included. Latent diffusion~\cite{rombach2022ldm} runs the process in the
latent space of a pretrained autoencoder and conditions it by cross-attention.

\textbf{Contrastive learning.} InfoNCE~\cite{oord2018cpc} and CLIP~\cite{radford2021clip}
align two encoders by asking each sample to identify its paired partner within a batch.

\pend{Person 3: add SAR-to-optical papers with diffusion and Transformers, and despeckling
work beyond Lee.}

% =====================================================================================
% Owner: Person 4
\section{Dataset and Split}

We use the Kaggle release ``Sentinel-1/2 image pairs segregated by
terrain''~\cite{requiemonk_dataset}, which follows the SEN1-2 format~\cite{schmitt2018sen12}.
It holds 16{,}000 co-registered pairs of 256$\times$256 tiles, 4{,}000 for each of four
terrains: agriculture, barren land, grassland and urban. The tiles are cut from 30 Sentinel
scenes. Each SAR tile is single-polarisation (VV), 8-bit. We read the grey values as a linear
mapping of $-25$ to $0$~dB, an assumption about how the dataset was exported.

\textbf{Scene split.} Whole scenes are held out (Table~\ref{tab:split}), and every terrain
has at least one scene in validation and in test. The validation scenes choose the best
epoch and the decoding settings. The test scenes are used only for the final scores. All
seven test scenes are fall scenes, because the dataset holds only one summer scene.

\begin{table}[t]
\centering
\caption{Scene-held-out split. Test tiles per terrain: agriculture 417, barren land 391,
grassland 449, urban 403.}
\label{tab:split}
\begin{tabular}{lrr}
\toprule
Part & Scenes & Pairs \\
\midrule
Train & 17 & 11{,}877 \\
Validation & 6 & 2{,}463 \\
Test & 7 & 1{,}660 \\
\midrule
Total & 30 & 16{,}000 \\
\bottomrule
\end{tabular}
\end{table}

% =====================================================================================
% Owner: Person 2 - expand from docs/METHOD_NOTES.md; all settings below are from bench_run.py.
\section{Proposed Method: PICMA}

PICMA (Physics-Informed Cross-Modal Attention latent diffusion) generates the optical tile
in the latent space of a frozen autoencoder. It is conditioned on the SAR tile through a
pretrained SAR encoder. Fig.~\ref{fig:arch} shows the full model.

\begin{figure}[t]
\centering
\fbox{\parbox[c][4.2cm][c]{0.92\columnwidth}{\centering\small
\pend{Person 5: architecture diagram}\\[2pt]
SAR $\rightarrow$ Lee filter $\rightarrow$ 3-ch physics input $\rightarrow$ ResNet-50
$\rightarrow$ (32$\times$32 spatial concat $|$ 320 tokens $\rightarrow$ cross-attention)
$\rightarrow$ latent U-Net $\times$ 50 DDIM steps $\rightarrow$ VAE decoder $\rightarrow$ RGB.
Side box: contrastive pretraining with optical ResNet-18. Loss box: $\epsilon$-MSE +
physics + FFL.}}
\caption{PICMA overview.}
\label{fig:arch}
\end{figure}

\subsection{Speckle Model and Physics Input}
SAR speckle is multiplicative: $I = R\,n$, where $R$ is the true reflectivity and $n$ has
mean 1 and variance $1/\mathrm{ENL}$. For Sentinel-1 IW GRD we use $\mathrm{ENL}=4.4$, so pure
speckle has coefficient of variation $C_u = 1/\sqrt{\mathrm{ENL}} \approx 0.48$. In a
7$\times$7 window on linear intensity, the Lee filter~\cite{lee1980speckle} computes the local
mean $m$, the local variation $C_i = \sqrt{v}/m$ and the weight
\begin{equation}
w = \operatorname{clip}_{[0,1]}\!\left(\frac{1 - C_u^2/C_i^2}{1 + C_u^2}\right),\qquad
I_f = m + w\,(I - m).
\end{equation}
Where $w \approx 0$ the window holds only speckle and is smoothed. Where $w \rightarrow 1$
the window holds real structure, such as an edge or a building, which is kept. Instead of
repeating the grey tile three times, the encoder receives three physics channels: raw SAR
(dB), Lee-filtered SAR (dB) and $C_i$ scaled by $\sqrt{\mathrm{ENL}}/4$.

\subsection{Physics Loss}
On the decoded prediction $(\hat L, \hat{ab})$ we add two terms:
\begin{equation}
\mathcal{L}_{\mathrm{phys}} = \overline{(1-w)\,\lVert\nabla \hat{ab}\rVert}
 + \overline{w\,\bigl(1 - \lvert\cos\angle(\nabla \hat L, \nabla I_f)\rvert\bigr)}.
\end{equation}
The first term keeps colour flat where the radar sees a homogeneous surface. The second term
aligns brightness edges with the despeckled radar edges. The absolute cosine allows the
contrast to flip between the sensors; a metal roof, for example, is bright in SAR but can be
dark in optical. The weight is 0.05.

\subsection{Cross-Modal Contrastive Pretraining}
Before diffusion training, the SAR encoder (ResNet-50, ImageNet initialisation) is aligned
with an optical encoder (ResNet-18, ImageNet initialisation) on the paired training tiles.
Each branch applies global average pooling and a two-layer MLP to a 128-d, L2-normalised
embedding. A symmetric InfoNCE loss~\cite{oord2018cpc,radford2021clip} (temperature 0.1, batch
64) asks each SAR tile to identify its own optical tile among the 64 tiles of the batch. This step runs
for 0.4~h at learning rate $10^{-4}$. Only the SAR encoder is kept.

\subsection{Latent Diffusion with Cross-Modal Attention}
A frozen Stable Diffusion VAE (\texttt{sd-vae-ft-mse})~\cite{rombach2022ldm} maps a
256$\times$256$\times$3 optical tile to a 32$\times$32$\times$4 latent (scale 0.18215). The
latent has 48 times fewer values than the tile. The forward process is DDPM~\cite{ho2020ddpm}
with $T=1000$ steps and a linear $\beta$ schedule from $10^{-4}$ to $0.02$. The denoiser predicts the noise
$\epsilon$ and is trained with $\mathcal{L}_\epsilon = \lVert\epsilon - \hat\epsilon\rVert_2^2$.

The denoiser is a latent U-Net with three levels (128/256/512 channels) and two residual
blocks per level. It has self-attention and cross-attention at 16$\times$16 and 8$\times$8. The SAR encoder
conditions it along two paths. In the \emph{spatial path}, layer2 features (32$\times$32,
512 channels) pass through a 1$\times$1 convolution to 64 channels and are concatenated with
the noisy latent, which gives pixel-aligned guidance. In the \emph{token path}, layer3
(16$\times$16) and layer4 (8$\times$8) give 320 tokens. These are projected to 512 dimensions and
receive a learned position embedding. In cross-attention~\cite{vaswani2017attention}, the
latent features are the queries and the SAR tokens are the keys and values. Every latent
position can therefore draw context from the whole tile.

\subsection{Auxiliary Losses on the Decoded Image}
Every fourth step, four tiles are noised to a low level ($t < 300$). The model predicts $x_0$
from $\hat\epsilon$, and the result is decoded with the VAE and converted to Lab. Two losses
are added on this decoded image: the physics loss (weight 0.05) and the focal frequency
loss~\cite{jiang2021ffl} (weight 1.0) between the 2-D spectra of the prediction and the
target. The FFL counters the blur that pixel losses cause. We use low $t$ only, because only
there is the decoded $x_0$ a meaningful image.

\subsection{Sampling}
At test time we use DDIM~\cite{song2021ddim} with 50 steps and clip the predicted $x_0$ to
the latent range. We draw $K=4$ samples per tile from different start noise. We report one
sample, which is a realistic output, and the mean of the $K$ samples, which is blurrier but
closer in pixel terms.

% =====================================================================================
% Owner: Person 1
\section{Experimental Setup}

\subsection{Compared Methods}
\begin{enumerate}
\item \textbf{Histogram match + mean colour.} L is the SAR intensity mapped onto the
  training L distribution by histogram matching (a lookup table from 2{,}000 training
  tiles). a/b is the mean training colour.
\item \textbf{Welsh}~\cite{welsh2002transferring}, adapted to SAR as a statistics lookup.
  Each pixel is described by its Lee-filtered intensity (32 bins) and its 5$\times$5 local
  standard deviation (8 bins). It receives the mean Lab colour (L, a and b) of the training
  pixels in the same cell, estimated from 2{,}000 training tiles. Cells with 20 or fewer
  pixels fall back to the global mean.
\item \textbf{U-Net R50}: our ResNet-50 U-Net colourizer, with a Zhang colour-bin
  head~\cite{zhang2016colorful} and a tanh L head. Its input is the SAR tile with Gaussian blurs at
  $\sigma=1$ and $2.5$. The encoder is fine-tuned at 0.05$\times$ the learning rate for 15
  epochs. We score it with 8-fold test-time augmentation.
\item \textbf{Pix2Pix}~\cite{isola2017pix2pix} \pend{run pending}.
\item \textbf{CycleGAN}~\cite{zhu2017cyclegan} \pend{run pending}.
\item \textbf{TransUNet}~\cite{chen2021transunet} \pend{teammate run, results not received}.
\item \textbf{Wavelet U-Net + FFL}~\cite{liu2018mwcnn,jiang2021ffl}
  \pend{teammate run, results not received}.
\item \textbf{DDPM-S}, a Palette-style~\cite{saharia2022palette} conditional pixel diffusion
  model. It generates L, a and b at 256~px, conditioned on the SAR tile at three
  despeckling levels. It uses a cosine schedule, EMA and DDIM sampling with 50 steps and
  $K=1$. It is a small model (5.6\,M parameters) trained for 2.5~h on 2$\times$T4
  (4{,}660 steps).
\item \textbf{PICMA (ours)}, Section~IV.
\end{enumerate}

\subsection{Training Details}
PICMA uses AdamW~\cite{loshchilov2019adamw} at learning rate $10^{-4}$ ($2\times10^{-5}$ for the SAR encoder), no weight
decay, batch 16, gradient clipping at 1.0 and mixed precision. The same random 90-degree
rotation and flip are applied to the SAR and the optical tile. An EMA of the weights (0.999)
is used at test time. Training runs on a fixed wall-clock budget on 2$\times$NVIDIA T4
(Kaggle), and the learning rate decays linearly over the second half of the budget.

The planned budget was 0.4~h contrastive pretraining plus 9.3~h diffusion training. That run
was stopped by the platform after 1.08~h of diffusion training (exit code $-9$, most likely
out of memory) and saved nothing. The PICMA results in this draft therefore come from a
backup run with a 1~h diffusion budget: 8 epochs and 5{,}239 steps, 1.52~h of training in
total. \pend{Replace if a longer run finishes.}

\subsection{Metrics}
All methods are scored by the same code on the 1{,}660 test tiles at 256~px. PSNR and
SSIM~\cite{wang2004ssim} are computed on RGB. a/b RMSE and L RMSE are computed in CIE Lab:
L is lightness, and a/b are the colour axes. LPIPS~\cite{zhang2018lpips} uses AlexNet, and
FID~\cite{heusel2017fid} uses 2048-d Inception features. FID on 1{,}660 images is biased
upward. It is valid for comparing our methods with each other, but not with published FID
values. DDPM-S was scored by its own evaluation script, which uses the same tiles and the
same metric definitions but does not compute FID.

% =====================================================================================
% Owner: Person 1
\section{Results}

\begin{table*}[t]
\centering
\caption{SAR-to-optical translation on the 7 held-out test scenes (1{,}660 tiles, 256~px).
Best value per column in bold. ``Train'' is GPU hours on 2$\times$T4. The mean-colour baseline
alone reaches an a/b RMSE of 9.13.}
\label{tab:results}
\begin{tabular}{llrrrrrrr}
\toprule
\# & Method & PSNR$\uparrow$ & SSIM$\uparrow$ & a/b RMSE$\downarrow$ & L RMSE$\downarrow$ &
LPIPS$\downarrow$ & FID$\downarrow$ & Train (h) \\
\midrule
1 & Hist.\ match + mean colour & 10.88 & 0.055 & \textbf{9.12} & 30.63 & 0.687 & 293.1 & 0 \\
2 & Welsh~\cite{welsh2002transferring} & \textbf{14.15} & 0.210 & 9.22 & \textbf{20.78} & \textbf{0.665} & 266.1 & 0 \\
3 & U-Net R50 (TTA) & 13.96 & \textbf{0.244} & 11.73 & 22.27 & 0.813 & 245.7 & \pend{h} \\
4 & Pix2Pix~\cite{isola2017pix2pix} & \multicolumn{7}{c}{\pend{pending}} \\
5 & CycleGAN~\cite{zhu2017cyclegan} & \multicolumn{7}{c}{\pend{pending}} \\
6 & TransUNet~\cite{chen2021transunet} & \multicolumn{7}{c}{\pend{pending}} \\
7 & Wavelet U-Net + FFL & \multicolumn{7}{c}{\pend{pending}} \\
8 & DDPM-S (pixel, $K{=}1$) & 9.97 & 0.051 & 20.09 & 31.68 & 0.802 & -- & 2.5 \\
\midrule
9 & PICMA (ours), sample & 7.77 & 0.089 & 32.80 & 40.09 & 0.820 & \textbf{244.9} & 1.52 \\
  & PICMA (ours), mean of $K{=}4$ & 11.10 & 0.186 & 19.70 & 29.09 & 0.783 & 258.9 & 1.52 \\
\bottomrule
\end{tabular}
\end{table*}

Table~\ref{tab:results} lists the finished runs. We note four findings.

\textbf{No method beats a constant colour on a/b error.} The best a/b RMSE (9.12) belongs to
histogram matching, whose a/b is just the mean training colour. Welsh transfer is close
(9.22). The learned models are further away: the U-Net reaches 11.73, which is 2.6 worse
than the mean colour. On scenes never seen in training, the SAR tile carries little
information about the true colour that a model can use. A model that commits to a colour
is penalised whenever the colour is wrong.

\textbf{Structure is predictable from SAR; colour is not.} Welsh transfer and the U-Net cut
the L RMSE from 30.63 (histogram match) to 20.78 and 22.27, and raise PSNR by about 3~dB.
Brightness and layout therefore follow the radar, while colour does not. The U-Net has the
best SSIM (0.244).

\textbf{Perception versus distortion.} PICMA's single samples have the lowest FID (244.9),
but the U-Net is only 0.8 behind (245.7). With 1{,}660 images this gap is within the noise of
FID, so we do not claim a win. PICMA is the worst method on every pixel metric. Averaging four
samples improves a/b RMSE from 32.80 to 19.70 and PSNR from 7.77 to 11.10~dB, but raises FID
to 258.9. This is the trade-off of Blau and Michaeli~\cite{blau2018perception}: the average
is more accurate per pixel and less realistic.

\textbf{Both diffusion models are undertrained.} The large spread between PICMA's samples
shows that, after 1~h, the denoiser does not yet follow the SAR condition closely. DDPM-S
behaves the same way at 2.5~h (a/b RMSE 20.09). Its authors report outputs that look like
satellite imagery but do not yet follow the SAR input. The PICMA budget was cut from the
planned 9.3~h to 1~h (Section~V-B), so these rows show the method at an early stage of
training. They do not show what the method can reach.

\begin{figure*}[t]
\centering
\fbox{\parbox[c][4cm][c]{0.9\textwidth}{\centering\small
\pend{Person 5: comparison grid - rows = test tiles, columns = SAR $|$ each method $|$ real
optical. Generate with \texttt{python scripts/bench_figure.py all\_results --out paper\_figs};
it needs each run's \texttt{shown.npz} and \texttt{metrics\_test.json}, which are not yet in
\texttt{all\_results/} (only logs were downloaded).}}}
\caption{Qualitative comparison on held-out test tiles.}
\label{fig:qual}
\end{figure*}

\subsection{Ablation}
\pend{Person 2: ablation set (same 2~h each: full PICMA, w/o physics, w/o contrastive
pretraining, w/o cross-attention). No ablation results received yet.}

% =====================================================================================
% Owner: Person 5
\section{Discussion and Limitations}

\begin{itemize}
\item \textbf{Compute.} All learned models trained for hours on T4 GPUs, not for days on large
  GPUs. The PICMA run with the full budget did not finish. Absolute numbers are therefore lower than those
  in published papers.
\item \textbf{Data.} We use VV polarisation only, and VH would help with vegetation. The dB scaling of the
  8-bit tiles is an assumption. All test scenes are fall scenes, so seasonal generalisation
  is untested.
\item \textbf{Metric.} RMSE-type metrics reward the mean colour on an ill-posed task.
  Perceptual metrics (LPIPS, FID) and visual inspection are needed alongside them. FID on
  1{,}660 images is noisy.
\end{itemize}
\pend{Person 5: failure cases from the comparison figure; where diffusion helps.}

% =====================================================================================
\section{Conclusion}
We presented PICMA, a latent diffusion model for SAR-to-optical translation. It uses the
speckle model in its input and its loss, aligns the SAR encoder with the optical domain by
contrastive pretraining, and reads SAR context through cross-attention. On a strict
scene-held-out benchmark, no method beats a constant mean colour on a/b error. Classical
statistics transfer leads on PSNR and LPIPS. PICMA, after only one hour of training, matches
the best FID but is the worst on pixel accuracy. Future work is a training run with the full
budget, the ablation study, VH polarisation, and multi-season test scenes.
\pend{rewrite once all rows are in.}

\bibliographystyle{IEEEtran}
\bibliography{refs}

\end{document}
